%! TeX root = ./main.tex
% vim: set spelllang=fr:
\section{mosaic --- Visualisation des vecteurs de Pareto}
\programdetails{https://github.com/matteodelabre/mosaic}{GPL v3.0}{Entre décembre~2023 et avril~2026, 200 lignes de code}

The cost space is a five-dimensional space in which each dimension corresponds to the possible cost values for each event type.
For a given input synteny and species tree, each point in this space represents the set of histories explicating the input and which are optimal with respect to the corresponding cost setting.
Notice that the event vector of any such optimal history is necessarily one of the Pareto-optimal output by Synesth.

Define the region of such a Pareto-optimal vector as the set of points (i.e., cost settings) which make the vector optimal with respect to the others.
Since there is a finite number of such vectors, there is also a finite number of such regions.
As shown by~\cite{LibeskindHadas2014ParetoOptimalPhylogenetic}, these regions are intersections of half-spaces and hence are convex polyhedrons.
These regions cover the entire cost space, but do not form a partition as neighboring polyhedra are not disjoint.

For visualization purposes, we reduce the five-dimensional space down to two dimensions, first by forcing the same cost for duplication and cut events and for transfer and transfer--cut events, and then by fixing the cost of one of the three remaining events to one (since the cost values are unit-less).
This enables a convenient representation of the cost space as a color-coded diagram, each color representing a given region.
For any given point of this diagram, its corresponding history set is the union of those of each region to which this point belongs.

